\begingroup%
\def\baselinestretch{1}%
\definecolor{D12Blue}{HTML}{3478B8}%
\definecolor{D12Teal}{HTML}{27877F}%
\definecolor{D12Amber}{HTML}{C58B28}%
\definecolor{D12Ink}{HTML}{243447}%
\definecolor{D12PaleBlue}{HTML}{F0F6FC}%
\definecolor{D12PaleTeal}{HTML}{EFF8F5}%
\begin{tikzpicture}[
  x=1mm, y=1mm,
  card/.style={draw, fill=white, rounded corners=2mm, line width=.75pt,
    inner xsep=2.5mm, inner ysep=2.5mm, outer sep=.4pt, align=flush center,
    text=D12Ink, font=\fontsize{9.5}{12.5}\selectfont},
  fast/.style={card, draw=D12Teal, minimum height=13mm},
  async/.style={card, draw=D12Amber, minimum width=72mm, minimum height=17mm},
  loop core/.style={card, draw=D12Blue, text width=42mm, minimum width=47mm, minimum height=29mm,
    font=\bfseries\fontsize{9.5}{12.5}\selectfont},
  group box/.style={draw, rounded corners=2.5mm, line width=.7pt,
    inner sep=0pt, outer sep=.4pt},
  group title/.style={inner sep=0pt, align=flush center, text=D12Ink,
    font=\bfseries\fontsize{10}{12.5}\selectfont},
  edge text/.style={fill=D12PaleTeal, inner xsep=1mm, inner ysep=.5mm, align=flush center,
    text=D12Ink, font=\fontsize{9}{12.5}\selectfont},
  flow/.style={draw=D12Teal, line width=.7pt, rounded corners=1.8mm,
    -{Latex[length=2.5mm,width=1.6mm]}},
  async flow/.style={flow, draw=D12Amber, dash pattern=on 2.6pt off 2pt},
  loop flow/.style={flow, draw=D12Blue, line width=.9pt,
    -{Latex[length=3mm,width=2mm]}}
]
\node[group title] at (0,12)
  {分级异步验证流水线：消除串行等待，严防 Reward Hacking 与隐式 Bug};
\node[fast, minimum width=51mm, minimum height=9mm] (cand) at (-41,0)
  {Coder 生成候选 Kernel};
\node[fast, minimum width=59mm] (g1) at (-41,-21)
  {Gate 1：轻量级同步门禁\\
   编译检查 · 静态校验 · 冒烟小 Shape};
\node[fast, minimum width=56mm] (next) at (45,-21)
  {立刻返回主循环继续下一轮迭代\\
   （主链路零阻塞）};
\node[async] (g2) at (-41,-49)
  {Gate 2：中层动态域与性能题\\
   跨字段动态 Shape · 掩码/空块边界 ·\\
   Geomean 测速};
\node[async] (g3) at (-41,-75)
  {Gate 3：重型上线准入压测\\
   全量盲测集 · Sanitizer 并发检测 ·\\
   防掉点监控};

\node[group title] at (0,-105)
  {自动化自进化 Loop 的头尾哲学：Data-Driven + Bench-Verified};
\node[loop core] (head) at (-53,-129)
  {头部：Data 数据集\\（第五章）\\
   Kernel Schema 与\\Schema Wiki\\
   提供高压缩先验，\\
   承担 Data-Driven 引擎};
\node[loop core] (mid) at (0,-129)
  {中部：Planner–Coder\\
   解耦闭环（第六章）\\
   设计空间树搜索 +\\
   受控工程恢复};
\node[loop core] (tail) at (53,-129)
  {尾部：Bench 测试集\\
   （第七章规划）\\
   分阶段分层级\\正确性与速度基准\\
   锚定物理真值，\\
   承担 Bench-Verified 准星};

\begin{scope}[on background layer]
  \node[group box, draw=D12Teal, fill=D12PaleTeal, fit={(-81,19) (81,-89)}] {};
  \node[group box, draw=D12Blue, fill=D12PaleBlue, fit={(-81,-98) (81,-150)}] {};
\end{scope}

\draw[flow] (cand) -- (g1);
\draw[flow] (g1) -- node[edge text] {快速通过 ·\\立即放行} (next);
\draw[async flow] (g1) -- node[edge text] {后台按重要度异步触发} (g2);
\draw[async flow] (g2) -- (g3);
\draw[async flow] (g3.east) -- (next.south |- g3.east)
  -- node[edge text] {发现隐式缺陷则告警回滚} (next.south);

\draw[loop flow] (head) -- (mid);
\draw[loop flow] (mid) -- (tail);
\end{tikzpicture}%
\endgroup%
